\subsection{Vacuum Configurations for Superstrings}

\textbf{Authors:} P. Candelas, Gary T. Horowitz, Andrew Strominger, Edward Witten

\paragraph{主な主張}
十次元 $E_8\times E_8$ 超弦・超重力を四次元 $N=1$ 超対称性を保って摂動的にコンパクト化するには内部六次元空間が $SU(3)$ holonomy を持つ Ricci-flat Kähler 多様体となり、standard embedding により四次元 $E_6$ 理論と $N_{\rm gen}-N_{\rm mirror}=\chi(K)/2$ の世代数が得られることを示した\cite{Candelas:1985en}。

\paragraph{新規性}
異常相殺条件と超対称性条件を Calabi--Yau 幾何、$E_8\to E_6$ のゲージ対称性破れ、Dirac index による世代数計数へ結び付け、自由作用する $\mathbb{Z}_5\times\mathbb{Z}_5$ 商による $\chi=-8$ の4世代模型を具体化した。

\paragraph{位置付け}
heterotic string phenomenology の出発点となった古典的研究であり、Calabi--Yau compactification、standard embedding、Euler characteristic による世代数計数、Wilson-line gauge breaking という後続の模型構築の基本枠組みを確立した。
